% Generated draft: requires mathematical and source review.
\documentclass[11pt]{article}
\usepackage[margin=25mm]{geometry}
\usepackage{fontspec}
\setmainfont{Latin Modern Roman}
\usepackage{amsmath,amssymb,mathtools,unicode-math}
\setmathfont{Latin Modern Math}
\usepackage{enumitem,needspace,xurl,hyperref}
\setlength{\parindent}{0pt}
\setlength{\parskip}{5pt}
\setlength{\emergencystretch}{3em}
\newcommand{\N}{\mathbb{N}}
\newcommand{\Z}{\mathbb{Z}}
\newcommand{\Q}{\mathbb{Q}}
\newcommand{\R}{\mathbb{R}}
\newcommand{\C}{\mathbb{C}}
\newcommand{\F}{\mathbb{F}}
\newcommand{\A}{\mathbb{A}}
\newcommand{\PP}{\mathbb{P}}
\newcommand{\E}{\mathbb{E}}
\newcommand{\eps}{\varepsilon}
\newcommand{\dd}{\,\mathrm{d}}
\begin{document}
\section*{UnsolvedMath ID: 10400039}
\textbf{Source problem code:} catalogue entry
\subsection{Definitions and mathematical statement}
A classical oriented knot is an oriented embedded circle in \(S^3\), up to isotopy. A long virtual knot is represented by a knot diagram of an oriented line, standard outside a compact set, with classical and virtual crossings, modulo the relevant Reidemeister and virtual detour moves. The precise virtual finite-type filtration is important and is not defined in the supplied statement.

The numbered claim [154], cross-referenced to Problem 3.9, is: every Vassiliev invariant \(v\) of classical knots admits an extension \(\widetilde v\) that is a finite-type invariant of long virtual knots. On a classical long knot \(k\), extension means
\[
 \widetilde v(k)=v(\operatorname{cl}(k)),
\]
where \(\operatorname{cl}(k)\) is its classical closure. The statement gives no coefficient group/ring, virtual finite-type convention, or bound comparing the two degrees. These are not silently supplied. In particular, classical-crossing and virtualization/semi-virtual formulations of virtual finite type must be distinguished when reviewing the claim.

The extraction then enters Section 2.8, ``Finite type invariants derived from local moves.'' This is surrounding context. It takes a commutative ring \(R\) with unit, the set \(\mathcal K\) of isotopy classes of oriented knots, and a local move \(m\) replacing a tangle inside a ball. If \(K\) admits \(m\) in \(d\) disjoint balls and \(K_S\) results from applying \(m\) in balls indexed by \(S\subset\{1,\ldots,d\}\), set
\[
 F_d(R\mathcal K,m)=\operatorname{span}_R
 \left\{\sum_{S\subset\{1,\ldots,d\}}(-1)^{|S|}K_S\right\}.
 \tag{13}
\]
The excerpt gives the descending filtration \(R\mathcal K=F_0\supset F_1\supset F_2\supset\cdots\) and calls an \(R\)-linear invariant \(v:R\mathcal K\to R\) finite type of \(m\)-degree \(d\) when \(v(F_{d+1})=0\); this condition bounds the degree and does not by itself assert exact degree. For the crossing-change move \(\times\), it writes \(F_d(R\mathcal K)=F_d(R\mathcal K,\times)\). It lists crossing change, double crossing change, \(\#\) move, pass move, \(\Delta\) move, doubled delta move, and \(n\)-gon move in Figure 10. The \(\Delta\) and \(n\)-gon strands may have any orientations matching across the move, and an \(n\)-gon side has \(n\) strands. The actual tangle pictures are missing. The contextual program asks to calculate the corresponding graded spaces; it does not provide a virtual-knot filtration for the numbered extension claim.
\subsection{Short English statement}
The main claim asks whether every classical Vassiliev invariant extends to a finite-type invariant of long virtual knots. The supplied registry calls this established by Goussarov--Polyak--Viro via Gauss-diagram formulas, while its title says conjecture and generic source evidence calls it open. These assessments and the unspecified virtual finite-type convention require source review. The following local-move discussion is contextual material.
\subsection{Sources}
\begin{itemize}
\item Local UnsolvedMath registry, permanent ID \(10400039\), exact source code \texttt{catalogue entry}, titled ``Conjecture 2.17'' followed by the extension claim [154]. The statement also says ``See also Problem 3.9.''
\item Ohtsuki, \emph{Problems on invariants of knots and 3-manifolds} (2002), Conjecture 2.17, PDF page 40, with extracted Section 2.8 context: \url{https://msp.org/gtm/2002/04/gtm-2002-04-024s.pdf}. The background distinguishes 231 printed bold numbered items from surrounding references.
\item Supplied reference [154]: Goussarov--Polyak--Viro, ``Finite type invariants of classical and virtual knots,'' \emph{Topology} 39 (2000), 1045--1068. The supplied research summary attributes the extension theorem and a Gauss-diagram construction to this reference and calls the virtual theory ``universal.''
\item Extraction pdf-font-structure, accessed 2026-07-29. Registry status solved and classification SOLVED-IN-LITERATURE coexist with a generic source-evidence NEEDS\_REVIEW note saying the item is open and requires a current primary-literature audit. No independent verification resolves that conflict here.
\item Dataset attribution: Ulam AI / UnsolvedMath Contributors, CC BY 4.0. Source rights NEEDS\_REVIEW: publicly accessible author/publisher copy, redistribution terms requiring review.
\end{itemize}
\end{document}
